\chapter*{Plan détaillé}
\addcontentsline{toc}{chapter}{Plan détaillé}

\section*{Fil conducteur proposé}
\addcontentsline{toc}{section}{Fil conducteur proposé}

J'aimerai partir sur l'hsitoire suivante : partir du problème biologique et méthodologique, montrer pourquoi un benchmark  était nécessaire, comparer les algorithmes existants, puis expliquer comment cette comparaison a motivé la création de scRAW, son optimisation par recherche d'hyperparamètres, le transfert de loss et les premiers tests pour le rendre inductif.

\section*{Question centrale possible}
\addcontentsline{toc}{section}{Question centrale possible}

\begin{center}
\fbox{\begin{minipage}{0.92\linewidth}
\textbf{Question directrice :} comment construire, évaluer et améliorer un pipeline de clustering scRNA-seq capable de mieux préserver les populations cellulaires rares ?
\end{minipage}}
\end{center}

Cette question peut servir à relier tous les blocs du rapport :
\begin{itemize}
  \item le benchmark répond au besoin de comparer proprement les méthodes ;
  \item l'état de l'art montre les forces et limites des approches existantes ;
  \item scRAW propose une réponse algorithmique centrée sur les populations rares ;
  \item la recherche d'hyperparamètres étudie la stabilité et les compromis ;
  \item l'inductif teste la capacité à passer d'une analyse fixé à un jeu de donnée à une prédiction sur de nouvelles données rendant l'application réelle encore plus concrète. Le transfert de loss peut permettre à transférer la loss de reconstruction pondérée à d'autres algorithmes de Deep Learning.
\end{itemize}

\section{Introduction}

\subsection*{Objectif de la section}
Introduire le contexte scRNA-seq, le besoin de clustering et le problème spécifique des populations rares.

\subsection*{Points à traiter}
\begin{itemize}
  \item Rappeler ce que sont les données scRNA-seq et pourquoi elles sont importantes.
  \item Expliquer le rôle du clustering dans l'identification de types cellulaires.
  \item Introduire les difficultés : grande dimension, bruit, sparsity, batch effects, déséquilibre des populations.
  \item Formuler le problème central : les populations rares sont mal capturées par beaucoup de pipelines. L'effet batch pertubre le clustering plus il est important
  \item Donner un exemple biologique concret de population rare : cellule immunitaire minoritaire, sous-type tumoral, population pathologique ou population spécifique d'un tissu.
  \item Expliquer pourquoi rater une population rare peut avoir un impact biologique : perte d'un signal d'intér\^et, confusion entre types proches, interprétation incomplète d'un tissu.
  \item Annoncer les contributions du stage : logiciel de benchmark, comparaison de l'état de l'art, création de scRAW, recherche d'hyperparamètres, tests inductifs, transfert de loss.
\end{itemize}

\subsection*{Autres éléments à inclure}
\begin{itemize}
  \item Mentionner que le clustering est souvent une étape amont pour l'annotation cellulaire et l'analyse biologique.
  \item Introduire l'élement fort de notre algorithme qui est l'équilibre entre performance globale et détection des petites populations.
  \item Parler de l'interet de la reproductibilité et de son application réelle dans l'analyse bioinformatique.
\end{itemize}

\section{Etat de l'art}

\subsection*{Plan possible}
\begin{enumerate}
  \item \textbf{Prétraitement des données scRNA-seq}
  \begin{itemize}
    \item filtrage cellules/gènes, normalisation, log-transformation, HVG, scaling ;
    \item importance de la reproductibilité du prétraitement.
    \item effets possibles du prétraitement sur les populations rares.
  \end{itemize}
  \item \textbf{Clustering classique}
  \begin{itemize}
    \item PCA + K-Means, PCA + Leiden ;
    \item avantages : simples, rapides, interprétables ;
    \item limites : sensibilité au bruit et aux populations rares.
  \end{itemize}
  \item \textbf{Approches par deep learning}
  \begin{itemize}
    \item autoencodeurs, scDeepCluster, scMAE, scNAME, (scCDCG, scVI-DCT ) ?? ;
    \item intér\^et : apprendre un espace latent plus adapté au clustering ;
    \item limites : co\^ut, hyperparamètres, reproductibilité, généralisation.
  \end{itemize}
  \item \textbf{Détection des populations rares}
  \begin{itemize}
    \item Le détection des classes rares est une problématique à laquelle des algorithmes ont tenté de répondre (Giniclust ? CellSIUS ? )
    \item expliquer pourquoi les classes rares contribuent peu à une loss moyenne ;
    \item introduire l'idée de pondération, de régularisation ou de post-traitement ;
  \end{itemize}
  \item \textbf{Effets de batch et généralisation}
  \begin{itemize}
    \item présenter le problème des variations techniques entre batch ;
    \item expliquer pourquoi un bon clustering transductif ne garantit pas une bonne prédiction sur de nouvelles cellules ;
    \item introduire l'intér\^et d'une évaluation inductive.
  \end{itemize}
  \item \textbf{Benchmarking des algorithmes}
  \begin{itemize}
    \item besoin d'une comparaison standardisée ;
    \item choix des métriques : ARI, NMI, ACC, métriques rares, batch metrics ;
    \item distinction entre protocole transductif et protocole inductif.
    \item importance du m\^eme prétraitement, des m\^emes seeds et des m\^emes exports pour comparer.
  \end{itemize}
  \item \textbf{Place de scRAW}
  \begin{itemize}
    \item motivation : rendre les cellules rares plus visibles pendant l'apprentissage ;
    \item positionnement : apprentissage spécifique et dirigé vers les populations rares et l'effet batch.
  \end{itemize}
\end{enumerate}

\subsection*{Figure ou tableau utile}
\begin{itemize}
  \item Tableau comparatif des algorithmes de l'état de l'art testés dans SCRBenchmark.
  \item Schéma simple : données scRNA-seq vers prétraitement, embedding, clustering, métriques.
  \item Figure conceptuelle montrant pourquoi les populations rares sont absorbées par les classes majoritaires.
\end{itemize}

\section{Matériel et méthodes}

\subsection*{Objectif de la section}
Dire précisément ce qui a été construit et comment les expériences ont été menées.

\subsection{Données et environnement expérimental}
\begin{itemize}
  \item Présenter les jeux de données utilisés : nom, espèce, tissu, nombre de cellules, nombre de batch.
  \item Mentionner l'environnement logiciel : Python, PyTorch, Scanpy, AnnData, Streamlit, GPU, SSH.
  \item Préciser les formats d'entrée et de sortie : \texttt{.h5ad}.
  \item Ajouter un tableau des datasets avec taille, nombre de classes, classes rares et protocole utilisé.
\end{itemize}

\subsection{Logiciel interactif de benchmarking}
\begin{itemize}
  \item Expliquer l'objectif de SCRBenchmark : centraliser chargement des données, sélection des algorithmes, exécution, visualisation, export.
  \item Décrire les deux modes : interface interactive et ligne de commande.
  \item Décrire la structure logique : registre d'algorithmes, gestion des données, métriques, visualisations.
  \item Expliquer pourquoi cet outil était nécessaire avant d'interpréter les résultats.
  \item Décrire les sorties sauvegardées : résultats bruts, tableaux agrégés, figures, configurations.
\end{itemize}

\subsection{Algorithmes de l'état de l'art testés}
\begin{itemize}
  \item Lister les algorithmes retenus et justifier leur présence.
  \item Séparer les baselines simples, les méthodes deep learning et les méthodes avec correction de batch.
  \item Préciser les paramètres communs et les conditions de comparaison.
\end{itemize}

\subsection{Description de scRAW}
\begin{itemize}
  \item Présenter le pipeline : prétraitement, autoencodeur, espace latent, pondération dynamique, DANN, Triplet Loss, clustering final.
  \item Expliquer la logique de pondération des cellules rares : cellules de petits pseudo-clusters ou zones peu denses ont plus de poids.
  \item Décrire les variantes testées.
  \item Expliquer les sorties produites par scRAW : embeddings, labels, poids cellulaires, modèle, figures.
  \item Distinguer la version stable, la version inductive et optimisée ?
  \item Ajouter un schéma du pipeline complet : données vers prétraitement, autoencodeur, pondération, clustering, métriques.
\end{itemize}

\subsection{Recherche d'hyperparamètres}
\begin{itemize}
  \item Décrire ce qui a été exploré : architecture, dimension latente, epochs, pondération, clustering, pertes, prétraitement latent.
  \item Préciser la stratégie 
  \item Expliquer que l'objectif n'est pas seulement le meilleur score global, mais le compromis global/rare.
  \item Mentionner les seeds, la répétabilité et la variabilité des résultats.
  \item Distinguer recherche large, ablations ciblées et consolidation finale.
  \item Décrire les critères d'arr\^et ou de sélection des configurations.
\end{itemize}

\subsection{Tests pour rendre scRAW inductif}
\begin{itemize}
  \item Distinguer transductif et inductif.
  \item Expliquer le protocole : train sur certains batch, gel du prétraitement et du modèle, prédiction sur batch non vus.
  \item Décrire les artefacts sauvegardés : modèle, état de prétraitement, référence centro\"ides, configuration.
  \item Expliquer pourquoi l'inductif est plus proche d'un usage réel : nouvelles cellules, nouveau batch, réanalyse rapide.
  \item Présenter les limites possibles : batch non représentatif, classes absentes du train.
\end{itemize}

\subsection{Métriques et analyses}
\begin{itemize}
  \item Expliquer les métriques globales : ARI, NMI, ACC.
  \item Expliquer les métriques centrées rares : RareACC, UltraRareACC et erreur par type cellulaire si présentées.
  \item Expliquer les métriques qualitatives : UMAP, matrices de confusion, visualisation des poids cellulaires.
  \item Discuter pourquoi une seule métrique peut \^etre trompeuse.
  \item Mentionner le traitement des clusters bruit/outliers par HDBSCAN.
\end{itemize}

\subsection{Reproductibilité et tra\c{c}abilité}
\begin{itemize}
  \item Sauvegarde des configurations exactes.
  \item Conservation des logs et des tableaux de résultats.
  \item Export des embeddings, labels prédits, poids et figures.
  \item Gestion des seeds et de l'environnement logiciel.
\end{itemize}

\section{Résultats et discussion}

\subsection*{Objectif de la section}
Raconter les résultats comme une progression : l'outil permet de comparer, la comparaison révèle des limites, scRAW répond à ces limites, puis l'optimisation et l'inductif testent sa solidité.

\subsection{Résultat 1 : logiciel de benchmarking}
\begin{itemize}
  \item Montrer ce que le logiciel permet de faire.
  \item Inclure capture ou schéma d'architecture.
  \item Discuter l'apport : facilité pour les biologistes non info, reproductibilité, comparaison, visualisation.
\end{itemize}

\subsection{Résultat 2 : comparaison des algorithmes existants}
\begin{itemize}
  \item Présenter un tableau global des scores.
  \item Identifier les forces/faiblesses des familles d'algorithmes.
  \item Mettre en avant les limites observées sur populations rares.
  \item Comparer les performances globales et les performances sur classes rares.
  \item Montrer les temps de calcul et la stabilité.
\end{itemize}

\subsection{Résultat 3 : performances et comportement de scRAW}
\begin{itemize}
  \item Montrer les scores principaux de scRAW par rapport aux autres meilleurs algorithmes.
  \item Discuter les gains, mais aussi les compromis.
  \item Montrer les UMAPs et UMAPs des poids ( figures JOBIM ?).
\end{itemize}

\subsection{Résultat 4 : recherche d'hyperparamètres}
\begin{itemize}
  \item Présenter les campagnes importantes.
  \item Montrer ce qui a été validé, invalidé ou jugé instable.
  \item Faire ressortir le compromis entre score global et détection des rares ?
  \item Discuter les configurations qui échouent.
  \item Montrer les effets du clustering final, de la loss, des poids. (afficher le plot Optuna ?)
\end{itemize}

\subsection{Résultat 5 : premiers tests inductifs}
\begin{itemize}
  \item Présenter le protocole train/test par batch.
  \item Montrer les résultats
  \item Discuter ce que ces tests disent de la généralisation de scRAW.
  \item Comparer l'inductif à une analyse transductive lorsque c'est possible.
\end{itemize}

\subsection{Résultat 6 : analyse qualitative des clusters}
\begin{itemize}
  \item Inclure UMAP, matrices de confusion ou erreurs par type cellulaire.
  \item Identifier les types cellulaires bien retrouvés et ceux qui restent difficiles. Se baser sur les marqueurs génétiques aussi ( voir ce qu'on fait ssCluBench pour l'analyse biologique)
  \item Discuter les clusters bruit ou non assignés.
\end{itemize}

\subsection{Résultat 7 : co\^ut de calcul et utilisabilité}
\begin{itemize}
  \item Comparer qualitativement les méthodes rapides, lourdes, stables ou difficiles à installer ( difficulté rencontré de version TensorFlow/pytorch par exemple).
  \item Mentionner le gain attendu du mode inductif pour réutiliser un modèle sans tout réntraîner.
\end{itemize}

\subsection{Discussion générale}
\begin{itemize}
  \item Ce qui fonctionne bien.
  \item Ce qui reste fragile.
  \item Limites : temps de calcul, sensibilité aux hyperparamètres, taille des datasets, validation biologique, origine de la vérité terrain.
  \item Discuter le compromis entre maximisation du score global et conservation des petites populations.
  \item Mentionner les risques de sur-optimisation sur un dataset de référence.
\end{itemize}

\section{Conclusion et perspectives}


\subsection*{Points à traiter}
\begin{itemize}
  \item Contribution logicielle : SCRBenchmark.
  \item Contribution méthodologique : comparaison standardisée des algorithmes.
  \item Contribution algorithmique : scRAW.
  \item Contribution expérimentale : recherche d'hyperparamètres et tests inductifs.
  \item Perspectives : à définir.
  \item Perspective logicielle : mieux packager scRAW, documenter l'installation, simplifier les relances.
  \item Perspective biologique : valider les populations rares avec des marqueurs ou une expertise externe.
\end{itemize}

\section{Figures et tableaux à préparer}

\begin{itemize}
  \item \textbf{Schéma général du pipeline de benchmarking} : données scRNA-seq, prétraitement, algorithmes, clustering, métriques et exports.
  \item \textbf{Capture ou schéma de SCRBenchmark} : montrer l'interface, les modes d'utilisation et les sorties produites.
  \item \textbf{Tableau des jeux de données} : nom, espèce, tissu, nombre de cellules, nombre de gènes, nombre de classes, classes rares et batchs.
  \item \textbf{Tableau des algorithmes testés} : famille de méthode, objectif, type d'entrée, prise en compte du batch, paramètres principaux et statut dans le benchmark.
  \item \textbf{Tableau des métriques} : Définir quelles métriques seront utilisées.
  \item \textbf{Schéma du pipeline scRAW} : prétraitement, autoencodeur, espace latent, pondération rare-aware, pertes utilisées, clustering final et exports.
  \item \textbf{Figure conceptuelle sur les populations rares} : illustrer pourquoi une petite population peut \^etre absorbée par des classes majoritaires.
  \item \textbf{Tableau principal des performances} : comparaison des algorithmes sur les métriques choisies 
  \item \textbf{UMAP comparatifs} : vérité terrain, meilleurs algorithmes de l'état de l'art et scRAW sur les datasets choisis.
  \item \textbf{Matrices de confusion} : montrer les populations bien retrouvées, les confusions entre types proches et les classes rares perdues.
  \item \textbf{Visualisation des poids cellulaires de scRAW} : UMAP coloré par poids o pour vérifier que les zones minoritaires sont mises en avant.
  \item \textbf{Tableau d'ablation} : Idem que Jobim.
  \item \textbf{Figure de recherche d'hyperparamètres} : résumé Optuna ou graphiques montrant les compromis entre performance globale, performance rare et stabilité.
  \item \textbf{Figure inductive} : schéma train/test par batch, puis comparaison entre résultats transductifs et inductifs et montrer que scRAW est meilleur.
  \item \textbf{Tableau du co\^ut de calcul} : performance en fonction du temps d'éxécution
  \item \textbf{Figure ou tableau d'analyse biologique} : À définir ce qu'on peut faire. 
\end{itemize}

